% Draft Introduction and Related Work (to be rewritten in unified voice by the writing agent).

\section{Introduction}

Somatic mutations drive the clonal evolution of cancer, and the resulting mixture of genetically distinct sub-populations provides both a window into tumour biology and a formidable computational challenge~\cite{nowell1976clonal,mcgranahan2017clonal}. Over the past decade, multi-region and longitudinal bulk DNA sequencing of tumours has repeatedly shown that intra-tumour heterogeneity is pervasive~\cite{gerlinger2012intratumor,nikzainal2012life,jamalhanjani2017tracking}. Multi-region sampling of non-small-cell lung cancer by the TRACERx consortium, for example, found that single-biopsy analysis would have missed up to 65\% of subclones and 76\% of subclonal mutations~\cite{jamalhanjani2017tracking}. As sequencing costs decline and study designs expand, multi-sample bulk datasets now routinely reach dozens of samples and dozens of subclones per patient, with some cohorts containing up to 90 samples and 26 subclones for a single patient~\cite{wintersinger2022pairtree}.

Reconstructing a tumour phylogeny from such data is a joint deconvolution and tree-inference problem: each bulk sample is modelled as a mixture of unknown clones, and the observed variant allele frequency (VAF) matrix $F$ must be factored simultaneously into a usage matrix $U$ and a clonal matrix $B$ while respecting the infinite-sites constraints of the perfect phylogeny model. A rich family of methods has been developed around this model, each with distinct loss functions, constraints and search strategies: AncesTree~\cite{elkebir2015reconstruction}, CITUP~\cite{malikic2015citup}, LICHeE~\cite{popic2015lichee}, PhyloWGS~\cite{deshwar2015phylowgs}, PASTRI~\cite{satas2017pastri}, CALDER~\cite{myers2019calder} and Pairtree~\cite{wintersinger2022pairtree}. Mutation-clustering front-ends such as PyClone~\cite{roth2014pyclone} and SciClone~\cite{miller2014sciclone} are often applied upstream to reduce the dimensionality of the problem.

A consistent observation across recent benchmarks is that these methods do not scale to the regime now defined by modern datasets. Pairtree~\cite{wintersinger2022pairtree} demonstrated that existing methods often fail beyond roughly ten subclones, and with the exception of the very recent method Orchard, all phylogenetic reconstruction methods operate on clusters of mutations rather than individual variant read counts, potentially discarding phylogenetic signal. Scalable inference that handles dozens of samples and hundreds of clones while retaining per-mutation resolution remains open.

In this work, we cast the VAF factorisation problem as a structured $\ell_1$-regression sub-problem over a fixed clonal tree nested within a combinatorial search over tree topologies, by analogy with minimum-evolution distance-based phylogenetics~\cite{desper2002fastme,lefort2015fastme,price2010fasttree}. Our main technical contribution is an $\mathcal{O}(mnd)$ algorithm for the $\ell_1$-VAF regression sub-problem that exploits the combinatorial structure of clonal matrices, together with an efficient recomputation procedure under subtree prune-and-regraft (SPR) moves. Building on this regression oracle, we develop fastBE, a simple deterministic search method for the $\ell_1$-VAF factorisation problem (VAFFP). On simulated data with up to 1000 clones and 100 samples, fastBE matches or exceeds the reconstruction accuracy of Pairtree and other competitors while running orders of magnitude faster. On 14 B-progenitor acute lymphoblastic leukaemia patients and two patient-derived colorectal cancer xenograft models~\cite{rehman2021diapause,ulintz2018polyclonal}, fastBE infers trees with fewer violations of the sum condition than Pairtree and reveals substantively different interpretations of intra-tumour heterogeneity.

Our contributions are: (i) a structured $\ell_1$-regression formulation of the VAF factorisation loss with an exact $\mathcal{O}(mnd)$ algorithm and SPR warm starts; (ii) fastBE, a deterministic beam-style search algorithm that couples this regression oracle to tree-space search; (iii) empirical evidence on simulated data showing mean wall-clock speedups of roughly two orders of magnitude over commercial linear programming solvers and strong scaling to 1000 clones; and (iv) biological analyses on B-ALL and colorectal PDX data showing that lower-loss phylogenies meaningfully differ in their heterogeneity estimates from those of the current state of the art.

\section{Background and Related Work}

\paragraph{Infinite-sites and perfect phylogenies.}
We follow previous work in restricting attention to somatic single-nucleotide variants in copy-neutral regions under an infinite-sites assumption~\cite{elkebir2015reconstruction,deshwar2015phylowgs,popic2015lichee,malikic2015citup,roth2014pyclone,satas2017pastri,myers2019calder,wintersinger2022pairtree}. Under this assumption, the evolutionary history of a tumour is a rooted phylogenetic tree $\mathcal{T}=(V,E)$ whose vertices correspond to clones and whose edges are labelled by the mutations gained along them. The $n$ clones form a vertex-labelled perfect phylogeny~\cite{elkebir2015reconstruction}, and the $n$-clonal matrix $B_{\mathcal{T}}$ records which mutations are present in each clone.

\paragraph{VAF factorisation.}
Each of the $m$ bulk samples is modelled as a convex combination of clones, giving $F \approx UB$ where $F$ is the observed frequency matrix and $U$ a right-stochastic usage matrix. Methods differ in the loss function that measures $F-UB$ and the additional constraints imposed: AncesTree~\cite{elkebir2015reconstruction} uses an $\ell_1$ loss with hard per-entry error constraints and solves via integer linear programming; CITUP~\cite{malikic2015citup} minimises a regularised $\ell_2$ loss by enumeration and quadratic integer programming over clusters of mutations; LICHeE~\cite{popic2015lichee} minimises squared violations of the sum condition via quadratic programming; CALDER~\cite{myers2019calder} extends AncesTree to longitudinal data using $\ell_0$ regularisation on $U$. Probabilistic methods such as PhyloWGS~\cite{deshwar2015phylowgs}, PASTRI~\cite{satas2017pastri} and Pairtree~\cite{wintersinger2022pairtree} model read counts directly and sample from a posterior over trees.

\paragraph{Distance-based phylogenetics.}
Minimum-evolution methods such as FastME~\cite{desper2002fastme,lefort2015fastme} and FastTree~\cite{price2010fasttree} have made large-scale distance-based phylogenetics practical by interleaving efficient branch-length regression with local topology moves (nearest-neighbour interchange and subtree prune-and-regraft). An analogous decomposition has been absent for tumour phylogenetics, where regression is typically solved either by integer linear programming or by general-purpose convex optimisation. Our regression algorithm provides the missing counterpart for the VAFFP.

\paragraph{Clustering front-ends and real datasets.}
Most methods rely on a mutation-clustering step such as PyClone~\cite{roth2014pyclone} or SciClone~\cite{miller2014sciclone} to reduce problem size. This clustering can discard phylogenetic signal; fastBE can operate either on mutation clusters or directly on individual mutations. The B-ALL dataset on which we benchmark contains fourteen patients with a median of 42 samples and up to 90 samples per patient~\cite{wintersinger2022pairtree}, and the colorectal xenograft models POP66 and CSC28 capture drug-resistant clonal dynamics~\cite{rehman2021diapause}. These datasets exemplify the scale at which current methods strain and motivate the algorithms developed here.
